% !TeX root = ../navier-stokes-manuscript.tex
% ============================================================
% 1.3 Gold Standard Proof
% ============================================================
\subsection{Gold Standard Proof}\label{sub:GoldStandardProof}

At this point the most direct way to solve the problem is visible.
Find one estimate produced by the initial datum, the fixed viscosity, and the Navier--Stokes equations that prevents the restart-grade Sobolev norm from escaping at a finite time.
For some fixed \(s>5/2\), the desired estimate has the form
\[
  \sup_{0\leq t<\min\{T,T_*\}}
  \lVert u(t)\rVert_{H^s(\mathbb R^3)}
  \leq
  C_s(u_0,\nu,T)
  <\infty
\]
for every finite \(T>0\).
This is the Gold Standard proof in one line.

The rest is already forced by what the bound controls.
Sobolev embedding keeps \(\nabla u\) bounded, the differentiated equations propagate each fixed higher \(H^m\) norm from the smooth datum, the pressure equation recovers the pressure rows, and the momentum equation recovers the time rows.
If \(T_*<\infty\), choose \(T=T_*\) and restart local theory from a time slice close enough to \(T_*\).
The same constant keeps the local lifespan from collapsing, so the restarted solution crosses the alleged last time and contradicts maximality.
Once the displayed estimate exists,
\[
  T_*=\infty.
\]

This route is so natural that the first attempt is almost unavoidable: differentiate the energy once and try to repeat the original argument at the gradient level.
Pair the momentum equation with \(-\Delta u\).
One integration by parts, together with incompressibility, turns the nonlinear transport into the cubic derivative integral
\[
  \frac12\frac{d}{dt}\lVert\nabla u\rVert_{L^2}^2
  +
  \nu\lVert\Delta u\rVert_{L^2}^2
  =
  -\int_{\mathbb R^3}
  \partial_k u_i\,\partial_i u_j\,\partial_k u_j\,dx.
\]
The cubic term is bounded by
\[
  \left|
  \int_{\mathbb R^3}
  \partial_k u_i\,\partial_i u_j\,\partial_k u_j\,dx
  \right|
  \leq
  C\lVert\nabla u\rVert_{L^3}^3
  \leq
  C\lVert\nabla u\rVert_{L^2}^{3/2}
  \lVert\Delta u\rVert_{L^2}^{3/2}.
\]
Young's inequality lets viscosity absorb the \(\Delta u\) factor and produces
\[
  \frac{d}{dt}\lVert\nabla u\rVert_{L^2}^2
  +
  \nu\lVert\Delta u\rVert_{L^2}^2
  \leq
  C\nu^{-3}\lVert\nabla u\rVert_{L^2}^6.
\]

The energy identity controls
\[
  \int_0^T\lVert\nabla u(t)\rVert_{L^2}^2\,dt,
\]
while the new row amplifies the same gradient through its sixth power.
A short high peak can have finite quadratic time area and still make the cubic derivative energy run away.
The inequality has exposed that analytic possibility without constructing such a fluid peak.
Its conclusion is the exact place where the first direct estimate loses control.

Every successful Gold Standard argument must reach into the same coupled equation and prevent that amplification at some continuation-grade level.
It has to use transport, pressure, incompressibility, and viscosity as they act on the original fluid, because those are the only mechanisms available to turn the local solution into an all-time one.
That missing estimate has resisted every direct attempt for almost a century.
